% Chapter 13: Rekeying, Closures, and Programs
% DEFINES: ch:rekeying-closures; sec:cipher-closure, sec:code-data-duality,
%   sec:rekeying, sec:programs, sec:ch13-notes;
%   def:cipher-closure, rem:designed-orbit, rem:rekeying-functor,
%   rem:control-flow
% REFERENCES (backward): rem:values-are-maps, sec:cipher-type, thm:orbit-bound,
%   thm:sum-impossibility, thm:composition
% REFERENCES (forward, part labels): part:frontiers

\chapter{Rekeying, Closures, and Programs}
\label{ch:rekeying-closures}

The two preceding chapters built cipher maps. This one builds \emph{with}
them: it shows that a single abstraction, the cipher closure, carries not
only values and functions but data structures, the secret itself, and whole
programs. The chapter folds three strands of the program, rekeying, cipher
data structures, and program realization, around that one idea, and where
each meets a limit, the limit is one the book has already proved. Control
flow runs into the sum-type impossibility of \Cref{thm:sum-impossibility};
sequential data leaks through the orbit closure of \Cref{thm:orbit-bound}.
The constructions are summarized from three companion drafts.

\section{The Cipher Closure}
\label{sec:cipher-closure}

Everything the untrusted machine handles in this framework has the same
shape: it carries a hidden secret and exposes an interface that operates on
bit strings. That shape has a name.

\begin{definition}[Cipher closure]
\label{def:cipher-closure}
A \emph{cipher closure} is a procedure that captures a trapdoor secret and
exposes an operation interface over bit strings, with the secret accessible
only through that interface. A cipher value (\Cref{rem:values-are-maps}), a
cipher map, and a cipher data structure are all cipher closures, differing
only in which interface they expose.
\end{definition}

The view is the closures-as-data-structures lens of structured programming,
applied to trapdoor computing: a closure is data that carries behavior, and
in this setting the captured datum is the secret. It unifies the objects of
the book. A cipher value is a closure whose interface is ``decode me'' (under
the trapdoor); a cipher map is a closure whose interface is ``apply me to a
token''; a cipher data structure, as the next section shows, is a closure
whose interface is a set of operations folded into one map. The secret is
captured in every case, and operated on, in the case of rekeying, by yet
another closure.

\section{Code Is Data, Data Is Code}
\label{sec:code-data-duality}

The closure makes precise a duality between cipher code and cipher data that
runs throughout the framework.

Data is code. A cipher data structure, a pair, a list, a map, is realized as
a \emph{single} cipher map. Its operation interface, the projections of a
pair, the head and tail of a list, is folded into a query tag: the untrusted
machine presents a tag naming the operation alongside the structure, and the
one cipher map dispatches on the tag to the right behavior. This is the
\emph{dispatch pattern}, and it turns a data structure into a function,
realizing data as code. Sequential structure is carried a second way, by a
\emph{designed orbit}: a list is encoded so that applying a hash-chain map
repeatedly walks from one element to the next, the structure living in the
shape of an orbit under that map.

\begin{remark}[Designed orbits leak length]
\label{rem:designed-orbit}
A structure encoded as an orbit pays the price \Cref{thm:orbit-bound} named.
The untrusted machine traverses a designed-orbit list by applying the
hash-chain map until the orbit closes, and the number of steps is the orbit
size, which is the list's length. So a designed-orbit data structure leaks
its length through exactly the orbit-closure channel of \Cref{ch:orbit-closure}:
the active adversary that walks the structure learns how big it is, even
though it never reads an element. Length is the residual the orbit bound
predicted.
\end{remark}

Code is data, in turn, because a cipher map is itself a point of a cipher
type, the cipher exponential $\cipher{A \to B}$ of \Cref{sec:cipher-type}. A
function is a value of a function type, so cipher code is cipher data, carried
like any other by a cipher map. The closure is the carrier that makes the
duality symmetric: values, functions, and data structures are all cipher
closures, and the cipher map is the universal vehicle for all of them.

\section{Rekeying}
\label{sec:rekeying}

If the secret is just another captured datum, it can be operated on, and the
operation that matters in practice is rotating it. Real systems change keys,
share data across trust boundaries, and retire old secrets; doing so naively
would mean decoding to plaintext and re-encoding under the new secret, which
defeats the point.

\begin{remark}[Rekeying is a cipher map]
\label{rem:rekeying-functor}
\emph{Cipher rekeying} maps a cipher value encoded under one secret to a
cipher value encoded under another, preserving the latent value, without
decoding. When the trusted machine knows both secrets it constructs the
rekeying transformation as a cipher map; once built, the \emph{untrusted}
machine applies it to any cipher value, rotating the key blindly. The
transformation is natural in the categorical sense: rekeying commutes with
the cipher operations, so a re-keyed value behaves under the algebra exactly
as the original did. Rekeying is then one instance of the closure principle,
a closure that transforms the captured secret of another closure.
\end{remark}

The information-theoretic cost of rekeying, and the conditions under which it
composes cleanly, are the rekeying paper's \cite{towell2026rekeying}; the
point for this part is structural: the secret is cipher data like any other,
and rotating it is a cipher map like any other.

\section{Programs as Cipher-Map Compositions}
\label{sec:programs}

The most ambitious use of the closure is to realize an entire program as a
composition of cipher maps the untrusted machine runs blindly. Two
difficulties shape what is possible, and both are limits the book has
already established.

The first is that ciphering is \emph{infectious}. Cutting a program at one
node, deciding to cipher that node's output, forces every node that consumes
the output to be ciphered too, since it now receives a cipher value, and
every sibling input that must be combined with it as well. The ciphered
region is therefore upward-closed toward the root, and a single cut can
cipher an entire single-rooted program. There is no local ciphering; the
decision propagates.

\begin{remark}[Control flow is sum-type dispatch]
\label{rem:control-flow}
The second difficulty is control flow, and it is the sum-type impossibility
of \Cref{thm:sum-impossibility} wearing operational clothes. Branching on a
cipher Boolean is untrusted dispatch on a sum type: the conditional must
choose a handler as a function of a hidden tag. By
\Cref{thm:sum-impossibility} it cannot both hide the tag and dispatch, so a
ciphered conditional faces exactly two options. It can \emph{expose} which
branch ran, letting the untrusted machine read the tag and run only the taken
branch, cheap but leaking the test; or it can \emph{hide} which branch ran,
evaluating both branches and combining them by a fused eliminator so the
untrusted machine never learns the tag, private but paying for both branches.
There is no third option, because the theorem forbids it. The cost of
oblivious control flow is the price of the sum-type wall.
\end{remark}

The full account of where to cut, how the ciphered region propagates, and how
control flow is compiled is the program-construction paper's
\cite{towell2026cipherprog}; here the realization closes the arc of the book.
The abstract composition theorem of \Cref{thm:composition} promised that
cipher maps chain; this chapter shows what chaining a whole program costs,
and that its two hard costs, infectious ciphering and oblivious branching,
are consequences of properties, totality and the sum-type impossibility, that
the earlier parts proved.

\section{Notes and Provenance}
\label{sec:ch13-notes}

The cipher closure and the code-data duality are the closures paper's
\cite{towell2026closures}; rekeying is \cite{towell2026rekeying}; the program
realization is \cite{towell2026cipherprog}. The chapter's contribution as a
chapter is to show the three as one: the cipher closure subsumes values,
maps, data structures, and the secret, and the limits the program
realization meets are not new obstructions but the orbit-closure leakage of
\Cref{ch:orbit-closure} and the sum-type impossibility of
\Cref{ch:algebraic-types}, now read operationally. With this chapter the
constructions of Part~V are complete: a hash-based batch construction, a
GF(2)-linear retrieval structure, and the closure machinery that carries data
structures, keys, and programs. What remains is the frontier, the problems
these constructions leave open, which is the book's final part.
